\documentclass[11pt]{amsart}
\usepackage{amsmath,amssymb,amsthm}
\usepackage[margin=1.1in]{geometry}
\usepackage{booktabs}
\usepackage[hidelinks]{hyperref}

\newtheorem{theorem}{Theorem}
\renewcommand{\thetheorem}{\Alph{theorem}}
\newtheorem{proposition}{Proposition}[section]
\newtheorem{lemma}[proposition]{Lemma}
\newtheorem{corollary}[proposition]{Corollary}
\theoremstyle{remark}
\newtheorem{remark}[proposition]{Remark}
\theoremstyle{definition}
\newtheorem{definition}[proposition]{Definition}

\newcommand{\Ga}{\Gamma}
\newcommand{\FK}{\Delta}
\newcommand{\tr}{\tau}
\newcommand{\dom}{\preccurlyeq}
\newcommand{\ltwo}{\ell^2}


\title{Product measures under invariant determinantal processes on groups}
\author{Yitzchak Shmalo}
\address{Einstein Institute of Mathematics, The Hebrew University of Jerusalem, Givat Ram, Jerusalem, Israel}
\subjclass[2020]{Primary 60G55; Secondary 60E15, 22D25, 43A07, 06F15}
\keywords{determinantal process, stochastic domination, Fuglede--Kadison determinant, amenable group, ordered group}
\date{}

\begin{document}

\begin{abstract}
Let $\Ga$ be a countable group and let $\mu_K$ be the determinantal measure on $\{0,1\}^{\Ga}$ given by a
$\Ga$-equivariant positive contraction $K$ on $\ltwo(\Ga)$. Lyons conjectured that $\mu_K$ stochastically dominates
the Bernoulli product measure of density equal to the Fuglede--Kadison determinant $\FK(K)$, is dominated by the product
measure of density $1-\FK(I-K)$, and that these bounds are optimal. We show that the density $\FK(K)$ can be exceeded
by a dominated product measure for some $K$ if and only if $\Ga$ is nonamenable, and likewise for the upper bound. We
also prove both inequalities for every left-orderable group, and more generally for groups with a subgroup $H$ such
that $\Ga/H$ carries a $\Ga$-invariant order and $H$ satisfies the lower inequality.
\end{abstract}

\maketitle

\section{Introduction}\label{sec:intro}

A determinantal process on a countable group $\Ga$ is given by an equivariant positive contraction $K$ on $\ltwo(\Ga)$,
and it is natural to ask how it compares with Bernoulli percolation. For $\Ga=\mathbb Z^d$, where $K$ is a Fourier
multiplier by a measurable $f$ with values in $[0,1]$, Lyons and Steif \cite{LS03} proved that $\mu_K$ stochastically
dominates the product measure whose density is the geometric mean $\exp\int\log f$, is dominated by the product measure
of density $1-\exp\int\log(1-f)$, and that neither density can be improved. Their argument reveals the sites one at a
time in lexicographic order and uses the Szeg\H{o}--Helson--Lowdenslager formula to bound the conditional probability of
each site given the sites already revealed.

On a general group the geometric mean is replaced by the Fuglede--Kadison determinant, and Lyons
\cite[Conj.~5.7]{LyICM} conjectured that for every equivariant positive contraction $K$,
\[
\pi_{\FK(K)}\dom\mu_K\dom\pi_{1-\FK(I-K)},
\]
and that these bounds are optimal, noting that the question is open even for finite groups. As in the theorem of
Lyons and Steif for $\mathbb Z^d$ \cite[Thm.~5.3]{LS03}, optimality means that for each $K$ no product measure of
density larger than $\FK(K)$ is dominated by $\mu_K$, and none of density smaller than $1-\FK(I-K)$ dominates it.

Here $\pi_p$ is the product measure of density $p$, $\mu_K$ is the determinantal measure of $K$, $\dom$ is
stochastic domination, and $\FK$ is the Fuglede--Kadison determinant; Section~\ref{sec:notation} gives the
definitions. Our first result shows that the optimality clause characterizes amenability; in particular it fails on
every nonamenable group.

\begin{theorem}\label{thm:A}
For a countable group $\Ga$ the following are equivalent.
\begin{enumerate}
\item $\Ga$ is amenable.
\item For every equivariant $0\le K\le I$ and every $p$, if $\pi_p\dom\mu_K$ then $p\le\FK(K)$.
\item For every equivariant $0\le K\le I$ and every $p$, if $\mu_K\dom\pi_p$ then $p\ge 1-\FK(I-K)$.
\end{enumerate}
\end{theorem}

On an amenable group
neither bound can be improved for any $K$; the upper bound on the density of a dominated product measure comes from
F{\o}lner approximation of the Fuglede--Kadison determinant \cite{LT14}, exactly as in the argument of Lyons and Steif.
On a nonamenable group there is always a kernel for which a product measure of density strictly larger than $\FK(K)$ is
dominated, and likewise at the upper end. The kernels are built from a symmetric random walk whose spectral radius is
small, which Kesten's theorem \cite{Kes59} provides; on the free group $F_2$ an explicit kernel is given in
Proposition~\ref{prop:F2}. Ray and Spinka \cite{RS} characterize amenability of graphs through stochastic domination
for the low-temperature Ising model and for the infinite cluster of Bernoulli percolation; Theorem~\ref{thm:A} is a
statement of the same kind for determinantal processes. Since $\pi_p=\mu_{pI}$, the domination $\pi_p\dom\mu_K$ holds
whenever $pI\le K$ \cite[Thm.~2.9]{LyICM}, and on finitely generated sofic groups Lyons and Thom \cite{LyT16} show
that such a monotone coupling can be chosen invariant; the inequalities studied here concern the threshold $\FK(K)$,
which is at least as large.

Our second result proves the two inequalities on a large class of groups: those having a subgroup $H$ for which
$\Ga/H$ carries a $\Ga$-invariant total order and on which the lower inequality holds.

\begin{theorem}\label{thm:B}
Let $H\le\Ga$ and suppose $\Ga/H$ carries a total order invariant under left multiplication by $\Ga$. Suppose that
every $H$-equivariant $0\le R\le I$ on $\ltwo(H)$ satisfies $\pi_{\FK(R)}\dom\mu_R$. Then every $\Ga$-equivariant
$0\le K\le I$ on $\ltwo(\Ga)$ satisfies
\[
\pi_{\FK(K)}\dom\mu_K\dom\pi_{1-\FK(I-K)}.
\]
\end{theorem}

With $H$ trivial
this gives both inequalities on every left-orderable group. In that case the proof is the coupling of Lyons and Steif,
run with the order of $\Ga$ in place of the lexicographic order, with Arveson's subdiagonal algebra of an ordered group
\cite{Arv67} in place of the Hardy space, and with the Jensen inequality for subdiagonal algebras
\cite{Lab05,BL} in place of the Szeg\H{o} formula. The determinant identity used there is a case of Austin's
Theorem~B \cite{Aus}: his theorem is stated for invariant random orders, and a left-invariant total order is the
point mass on one of them \cite[Sec.~4.1]{Aus}, where the mean in his formula has a single term. Theorem~\ref{thm:B}
joins this identity to a coupling that gives stochastic domination, and allows a nontrivial $H$. The coset form allows torsion in $H$; for instance the
inequalities hold on $C_2\times F_2$, where $C_2$ is the group of order two and $F_2$ the free group on two generators
(Corollary~\ref{cor:finite}).

Section~\ref{sec:notation} fixes conventions and proves a conditioning lemma used in both proofs. Sections~\ref{sec:A}
and~\ref{sec:B} prove Theorems~\ref{thm:A} and~\ref{thm:B}.

\section{Notation and conventions}\label{sec:notation}

Throughout, $\Ga$ is a countable group, $\lambda$ and $\rho$ are its left and right regular representations on
$\ltwo(\Ga)$, and $R(\Ga)$ is the von Neumann algebra generated by $\rho(\Ga)$, with trace
$\tr(T)=\langle Te_{e},e_{e}\rangle$; $T^*$ denotes the adjoint of $T$. A bounded operator $K$ on $\ltwo(\Ga)$ is \emph{equivariant} if it commutes with
$\lambda(\Ga)$, that is, with left translations; equivalently $K\in R(\Ga)$, $K$ acts by right convolution, and
$\langle Ke_y,e_x\rangle$ depends only on $x^{-1}y$.
With this convention $\rho(g)e_x=e_{xg^{-1}}$, so $e_x=\rho(x^{-1})e_e$ for every $x$. Results stated in the literature
for the opposite convention, with operators commuting with right translations, transfer to ours through the unitary
$e_x\mapsto e_{x^{-1}}$, which carries $R(\Ga)$ onto the left group von Neumann algebra $L(\Ga)$, preserves the trace,
and carries the compression of an operator to a finite set $F$ to the compression of its image to $F^{-1}$.

For a positive operator $T\in R(\Ga)$ the Fuglede--Kadison determinant is $\FK(T)=\exp\tr(\log T)$, with value $0$ when
$\tr(\log T)=-\infty$. For $0\le K\le I$ equivariant, $\mu_K$ is the determinantal probability measure on $\{0,1\}^{\Ga}$
with $\mu_K[\,F\subseteq\omega\,]=\det(K_F)$ for finite $F\subseteq\Ga$. Write $\pi_p$ for the product measure of density
$p$, and $\nu\dom\nu'$ for stochastic domination.

Domination of infinite-volume measures is always obtained as a limit of domination on finite windows; a window and an
order of revelation inside it are fixed before any coupling is built.

\subsection*{Schur complements and conditioning}
Let $E$ be countable, let $Q$ be a bounded positive operator on $\ltwo(E)$, and let $A\subseteq E$ be such that the
compression $Q_A$ of $Q$ to $\ltwo(A)$ satisfies $Q_A\ge\varepsilon I$ for some $\varepsilon>0$. Write
$Q^{\langle A\rangle}$ for the operator on $\ltwo(E\setminus A)$ given by
\[
\langle Q^{\langle A\rangle}f,f\rangle=\inf\bigl\{\langle Q(f-c),f-c\rangle:\ c\in\ltwo(A)\bigr\},\qquad
f\in\ltwo(E\setminus A),
\]
the Schur complement of $Q$ relative to $A$. Because $Q_A$ is invertible the infimum is attained at
$c=Q_A^{-1}Q_{A,E\setminus A}f$, so $Q^{\langle A\rangle}=Q_{E\setminus A}-Q_{E\setminus A,A}Q_A^{-1}Q_{A,E\setminus A}$ is a
bounded positive operator. It is monotone: $Q_1\le Q_2$ implies $Q_1^{\langle A\rangle}\le Q_2^{\langle A\rangle}$, since
the infimum is taken of pointwise ordered forms; and $Q^{\langle A\rangle}\le Q_{E\setminus A}$, by taking $c=0$.

\begin{lemma}[conditioning]\label{lem:cond}
Let $E$ be finite, let $\varepsilon I\le K\le(1-\varepsilon)I$ on $\ltwo(E)$ for some $\varepsilon>0$, and let
$A,B\subseteq E$ be disjoint. Conditioned on $A\subseteq\omega$ and $B\cap\omega=\emptyset$, the restriction of $\omega$ to
$E\setminus(A\cup B)$ is determinantal with a kernel $K'$ satisfying
\[
K'\ \ge\ \bigl(K^{\langle A\rangle}\bigr)_{E\setminus(A\cup B)}\ \ge\ K^{\langle A\cup B\rangle}.
\]
\end{lemma}

\begin{proof}
The bound $K\le(1-\varepsilon)I$ gives $\mu_K(B\cap\omega=\emptyset)=\det(I-K)_B>0$, so the first conditioning below is
defined. First condition on $B\cap\omega=\emptyset$. The complement $1-\omega$ has law $\mu_{I-K}$, and the event is
$\{B\subseteq1-\omega\}$. By \cite[Prop.~8.3]{Ly03}, which holds for every positive contraction (on a finite $E$ it is the block-determinant
identity $\det Q_{A\cup F}/\det Q_A=\det(Q^{\langle A\rangle})_F$ for $A\cap F=\emptyset$), the conditioned
complement on $E\setminus B$ is determinantal with kernel $\langle P^{\perp}e,P^{\perp}f\rangle_{I-K}$, where
$\langle u,v\rangle_{I-K}=\langle(I-K)u,v\rangle$ and $P^\perp$ is the projection orthogonal to $\ltwo(B)$ in this inner
product; that kernel is the Schur complement $(I-K)^{\langle B\rangle}$. So $\omega$ on $E\setminus B$ is determinantal
with kernel $K^{B}:=I-(I-K)^{\langle B\rangle}\ge I-(I-K)_{E\setminus B}=K_{E\setminus B}$. Next condition this process
on $A\subseteq\omega$, an event of probability $\det(K^B)_A>0$ because $K^B\ge K_{E\setminus B}\ge\varepsilon I$. By
\cite[Prop.~8.3]{Ly03} again, the result on $E\setminus(A\cup B)$ is determinantal with kernel $K'=(K^B)^{\langle
A\rangle}$. Monotonicity of the Schur complement and $K^B\ge K_{E\setminus B}$ give
$K'\ge(K_{E\setminus B})^{\langle A\rangle}=(K^{\langle A\rangle})_{E\setminus(A\cup B)}$, the last equality because the
infimum runs over $c\in\ltwo(A)$ and $f$ is supported off $B$. Enlarging the set of admissible $c$ from $\ltwo(A)$ to
$\ltwo(A\cup B)$ can only lower the infimum, which gives the second inequality.
\end{proof}

\section{Theorem A: amenability}\label{sec:A}

\begin{proof}[Proof that (1) implies (2)]
Suppose $\pi_p\dom\mu_K$. For a finite set $F\subseteq\Ga$ the event $\{F\subseteq\omega\}$ is increasing, so
\[
p^{|F|}=\pi_p[F\subseteq\omega]\le\mu_K[F\subseteq\omega]=\det K_F ,
\]
the argument of Lyons and Steif \cite{LS03} for $\Ga=\mathbb Z^d$. The map $e_x\mapsto e_{x^{-1}}$ is a unitary of
$\ltwo(\Ga)$ that carries $R(\Ga)$ onto $L(\Ga)$, preserves $\tr$, and carries the compression $K_F$ to the compression
of the image of $K$ to $F^{-1}$; so $\FK$ and the set of numbers $(\det K_F)^{1/|F|}$, $F$ finite, are unchanged.
For an amenable group and any positive element, Li and Thom \cite[Thm.~1.4]{LT14} prove that the Fuglede--Kadison
determinant equals $\inf_F(\det K_F)^{1/|F|}$ over finite $F$ (and the limit along F{\o}lner sets); no invertibility
is assumed, so the case $\FK(K)=0$ is included. Letting $F$ run along a F{\o}lner sequence gives $p\le\FK(K)$.
\end{proof}

\begin{proof}[Proof that (2) fails on nonamenable groups]
Let $\Ga$ be nonamenable. By the F{\o}lner criterion some finitely generated subgroup $\Ga_0$ is nonamenable, and by
Kesten's theorem \cite{Kes59} the uniform probability measure $\nu_0$ on a finite symmetric generating set of $\Ga_0$
has a convolution operator $P_{\nu_0}$ of norm $u_0<1$. The powers $\nu_0^{*n}$ have norms $u_0^n$, and their mass
$b_n=\nu_0^{*n}(e)=\langle P_{\nu_0}^n e_e,e_e\rangle$ at the identity is at most $u_0^n$. Removing that mass and
renormalising gives a symmetric probability measure $\nu=(\nu_0^{*n}-b_n\delta_e)/(1-b_n)$ with $\nu(e)=0$ and
$q:=\|P_\nu\|\le(u_0^n+b_n)/(1-b_n)$; fix $n$ with $q<1/2$. Regarded on $\Ga$, the operator $P_\nu$ splits along the
cosets of $\Ga_0$ into copies of itself, so its norm is still $q$.

Define $M$ on $\ltwo(\Ga)$ by $\langle Me_y,e_x\rangle=\sqrt{\nu(x^{-1}y)}$. It is equivariant, real, self-adjoint and of
finite range, with $\tr(M)=0$ and $\tr(M^2)=\sum_g\nu(g)=1$. For a nonempty finite set $F$,
\[
\sum_{x,y\in F}|M_{xy}|^2=\langle 1_F,P_\nu 1_F\rangle\le q|F|.
\]
Hence some $x\in F$ has $\sum_{y\in F}|M_{xy}|^2\le q$. Place such an $x$ last, remove it, and repeat. In the
resulting order $x_1,\dots,x_m$ of $F$ each $x_j$ satisfies $\sum_{i<j}|M_{x_jx_i}|^2\le q$.

Let $K=\tfrac12I+tM$ with $|t|\,\|M\|<\tfrac12$, so that $K$ and $I-K$ are bounded below by
$\kappa:=\tfrac12-|t|\,\|M\|$. The restriction of $\mu_K$ to $F$ is $\mu_{K_F}$, and $\kappa I\le K_F\le(1-\kappa)I$.
Reveal $\omega$ on $F$ in the order above. Apply Lemma~\ref{lem:cond} to the compression $K_E$ of $K_F$ to $E=\{x_1,\dots,x_j\}$ with
$\varepsilon=\kappa$, with $A$ and $B$ the revealed points of $\{x_1,\dots,x_{j-1}\}$ that lie in and out of $\omega$; its
second inequality, with $A\cup B=\{x_1,\dots,x_{j-1}\}$, shows that whatever has been revealed, the conditional
probability that $x_j\in\omega$ is at least the Schur complement
\[
K_{x_jx_j}-k^{*}(K_{\{x_1,\dots,x_{j-1}\}})^{-1}k\ \ge\ \tfrac12-\frac{t^2q}{\kappa}=:p(t),
\qquad k=(K_{x_ix_j})_{i<j},
\]
since $\|k\|^2=t^2\sum_{i<j}|M_{x_ix_j}|^2\le t^2q$ and the compression is bounded below by $\kappa$. Revealing
$\omega$ site by site against independent uniform variables then couples $\pi_{p(t)}$ below $\mu_K$ on $F$
(this is Holley's lemma in the form \cite[Lemma~5.2]{LS03}). Domination is preserved under weak limits, since by
\cite{Str65} it amounts to inequalities between the probabilities of increasing cylinder events; letting $F$ increase
gives $\pi_{p(t)}\dom\mu_K$.

On the other hand $\log(\tfrac12 I+tM)=\log\tfrac12\,I+\sum_{k\ge1}(-1)^{k+1}(2tM)^k/k$, so
$\FK(K)=\tfrac12\exp\bigl(-2t^2+O(|t|^3)\bigr)=\tfrac12-t^2+O(|t|^3)$, where the implied constant depends only on
$\|M\|$. Since $p(t)=\tfrac12-2q t^2+O(|t|^3)$ and $2q<1$, we get $p(t)>\FK(K)$ for all sufficiently small
$t\neq0$, so (2) fails. The same estimate applied to $I-K=\tfrac12I-tM$ shows that (3) fails.
\end{proof}

\begin{proof}[Proof that (2) and (3) are equivalent]
If $\omega$ has law $\mu_K$ then $1-\omega$ has law $\mu_{I-K}$ \cite{Ly03}, and $\omega\mapsto1-\omega$ reverses the
pointwise order. Hence $\mu_K\dom\pi_p$ if and only if $\pi_{1-p}\dom\mu_{I-K}$. Since $K\mapsto I-K$ is a
bijection of the equivariant positive contractions, statement (3) for $K$ is statement (2) for $I-K$ with $1-p$ in place
of $p$.
\end{proof}

\begin{proposition}\label{prop:F2}
Let $\Ga=F_2$ with free generators $a,b$, let $Y=\rho(a)+\rho(a^{-1})+\rho(b)+\rho(b^{-1})$ and
$K=\tfrac12 I+\tfrac1{32}Y$. Then $\pi_{191/384}\dom\mu_K$ and $\mu_K\dom\pi_{193/384}$, while
$\FK(K)<64/129<191/384$ and $1-\FK(I-K)>65/129>193/384$.
\end{proposition}

\begin{proof}
The matrix of $Y$ is the adjacency matrix of the Cayley graph of $F_2$ for $\{a^{\pm1},b^{\pm1}\}$, the $4$-regular
tree $T$, so $\|Y\|\le4$ and $\tfrac38I\le K\le\tfrac58I$. Fix a ball $F$ of $T$ about $e$ and order it by distance
from $e$, breaking ties arbitrarily. Each $x\ne e$ in $F$ has exactly one neighbour before it, namely its parent, since
its other neighbours are farther from $e$ and vertices at its own distance are not adjacent to it. Lemma~\ref{lem:cond},
applied as in the proof of Theorem~\ref{thm:A} with $\kappa=\tfrac38$ and $\|k\|^2=(1/32)^2$, shows that given any
revealed values on the earlier points, $x$ lies in $\omega$ with conditional probability at least
\[
\tfrac12-\frac{1/1024}{3/8}=\tfrac12-\tfrac1{384}=\tfrac{191}{384},
\]
and at $x=e$, which comes first, the probability is $K_{ee}=\tfrac12$. The same lemma for $I-K=\tfrac12I-\tfrac1{32}Y$, whose law is that of
$1-\omega$, shows that $x$ is absent with conditional probability at least $\tfrac{191}{384}$. The sequential coupling
\cite[Lemma~5.2]{LS03} and the limit over balls give $\pi_{191/384}\dom\mu_K\dom\pi_{193/384}$.

For the determinant, $\|Y/16\|\le\tfrac14$, so
$\tr\log(I+Y/16)=\sum_{k\ge1}(-1)^{k+1}\tr(Y^k)/(k\,16^k)$. Here $\tr(Y^k)$ counts closed walks of length $k$ at $e$
in $T$; it vanishes for odd $k$, since $T$ is bipartite, and every even term enters with a minus sign. Keeping only
$k=2$, with $\tr(Y^2)=4$,
\[
\FK(K)=\tfrac12\exp\tr\log(I+Y/16)\le\tfrac12e^{-1/128}<\tfrac12\Bigl(1-\tfrac1{128}+\tfrac1{2\cdot128^2}\Bigr)
<\tfrac12\cdot\tfrac{128}{129}=\tfrac{64}{129}.
\]
The same series for $I-K=\tfrac12(I-Y/16)$ has the same even terms and no odd ones, so $\FK(I-K)<\tfrac{64}{129}$ as
well. Finally $191\cdot129-64\cdot384=63>0$ and $65\cdot384-193\cdot129=63>0$.
\end{proof}

\section{Theorem B: groups ordered modulo a subgroup}\label{sec:B}

\begin{corollary}\label{cor:lo}
Lyons's inequalities hold on every left-orderable group.
\end{corollary}

\begin{corollary}\label{cor:finite}
Let $H$ be a normal subgroup of $\Ga$ with $\Ga/H$ left-orderable. If $H$ has order $2$, or more generally if $H$
satisfies the hypothesis of Theorem~\ref{thm:B}, then Lyons's inequalities hold on $\Ga$; for instance on $C_2\times F_2$.
\end{corollary}

\begin{remark}
For $\Ga/H=\mathbb Z^d$ and $H$ finite, an equivariant kernel is a $\mathbb Z^d$-stationary matrix-valued kernel, the
setting of Cyr \cite{Cyr21}, whose study of such processes includes a notion of full domination. The bounds of
Theorem~\ref{thm:B} concern stochastic domination and the geometric mean $\FK(K)$.
\end{remark}

\begin{proof}[Proof of Theorem~\ref{thm:B}]
Write $<$ for the order on $\Ga/H$, and put
\[
U_+=\{x\in\Ga:\ xH>H\},\qquad S=\{s\in\Ga:\ s^{-1}H>H\}=U_+^{-1}.
\]

\emph{Step 1: a subdiagonal algebra.} Left invariance of the order gives $SS\subseteq S$: if $s_1,s_2\in S$, applying
$s_2^{-1}$ to $s_1^{-1}H>H$ gives $(s_1s_2)^{-1}H>s_2^{-1}H>H$. It also gives $HS\subseteq S$ and $SH\subseteq S$, since
for $h\in H$ and $s\in S$ we have $(hs)^{-1}H=s^{-1}H$ and $(sh)^{-1}H=h^{-1}s^{-1}H>h^{-1}H=H$. For every $x\in\Ga$ exactly one of $xH=H$, $xH>H$,
$xH<H$ holds, and applying $x^{-1}$ to $xH<H$ gives $x^{-1}H>H$; so $\Ga$ is the disjoint union of $H$, $S$ and
$S^{-1}$. Let $\mathcal A$ be the weak$^*$-closed linear span of $\{\rho(g):g\in H\cup S\}$ in $R(\Ga)$, let
$\mathcal A_0$ be that of $\{\rho(s):s\in S\}$, and let $\Phi$ be the trace-preserving conditional expectation of
$R(\Ga)$ onto $R(H)$, which keeps the coefficients on $H$ and drops the rest. By the relations just shown, $\mathcal A$
is an algebra, $\mathcal A_0$ is an ideal in it with $\Phi(\mathcal A_0)=0$, $\Phi$ is multiplicative on $\mathcal A$,
$\mathcal A\cap\mathcal A^*=R(H)$, and $\mathcal A+\mathcal A^*$ contains every $\rho(g)$ and so is weak$^*$ dense.
(Multiplicativity of $\Phi$ passes from the span to its weak$^*$ closure because $\Phi$ is normal and multiplication is
separately weak$^*$ continuous; $\mathcal A\cap\mathcal A^*=R(H)$ because an element of both has its Fourier coefficients
supported in $(H\cup S)\cap(H\cup S^{-1})=H$.) Thus $\mathcal A$ is a weak$^*$-closed unital subalgebra of the finite
von Neumann algebra $R(\Ga)$ satisfying Arveson's conditions \cite{Arv67}. These are the conditions in the definition
of a (finite maximal) subdiagonal algebra in \cite[Sec.~1]{BL}. For such algebras Jensen's inequality $\FK(x)\ge\FK(\Phi(x))$ holds for every $x$ in the algebra,
invertible or not \cite[Cor.~3.3]{BL}. Hence
\begin{equation}\label{eq:jensen}
\FK(x)\ge\FK(\Phi(x))\qquad(x\in\mathcal A).
\end{equation}

\emph{Step 2: the Schur complement on $H$.} Fix $\delta\in(0,\tfrac12)$ and put $K_\delta=(1-2\delta)K+\delta I$, so
that $\delta I\le K_\delta\le(1-\delta)I$. Define $R$ on $\ltwo(H)$ by
\[
\langle Rf,f\rangle=\inf\bigl\{\langle K_\delta(f-b),f-b\rangle:\ b\in\ltwo(U_+)\bigr\},
\]
that is, $R=K_\delta^{\langle U_+\rangle}$ compressed to $\ltwo(H)$; since the compression of $K_\delta$ to
$\ltwo(U_+)$ is at least $\delta$, this is the usual operator Schur complement, and $\delta I\le R\le I$. For $h\in H$,
$\lambda(h)$ maps $H$ and $U_+$ onto themselves (left invariance again) and commutes with $K_\delta$, so $R$ commutes
with $\lambda(H)$ on $\ltwo(H)$, that is, $R\in R(H)$. Through the inclusion $R(H)\subseteq R(\Ga)$ that sends
$\rho(h)$ to $\rho(h)$ we regard $R$ also as an operator on $\ltwo(\Ga)$; the trace of $R(\Ga)$ restricts to that of
$R(H)$, so $\FK(R)$ is the same in both algebras, and $\FK(R)$ in the hypothesis of Theorem~\ref{thm:B} is taken in
this sense.

\emph{Step 3: $\FK(R)\ge\FK(K_\delta)$.} Let $d\in R(H)$. Then $de_e\in\ltwo(H)$, finitely supported vectors are dense in
$\ltwo(U_+)$, and a finitely supported $b\in\ltwo(U_+)$ has the form $b=ae_e$ with $a$ in the linear span of
$\{\rho(s):s\in S\}$, because $e_x=\rho(x^{-1})e_e$. Since $\langle Te_e,e_e\rangle=\tr(T)$ for $T\in R(\Ga)$,
\[
\tr(d^*Rd)=\langle Rde_e,de_e\rangle=\inf_a\ \tr\bigl((d-a)^*K_\delta(d-a)\bigr),
\]
the infimum over that span. For each such $a$, $x=d-a$ lies in $\mathcal A$ with $\Phi(x)=d$, so by the arithmetic--
geometric inequality $\tr(y)\ge\FK(y)$ for $y\ge0$ (that is, $\log\tr(y)\ge\tr(\log y)$, the classical Jensen inequality
for the concave function $\log$ and the spectral measure of $y$ in the state $\tr$), the identities $\FK(ab)=\FK(a)\FK(b)$ and $\FK(a^*)=\FK(a)$, which
hold for arbitrary, not necessarily invertible, elements of a finite von Neumann algebra \cite[Thm.~2.1]{BL} (with
proofs credited there to Haagerup and Schultz), and \eqref{eq:jensen},
\[
\tr(x^*K_\delta x)\ \ge\ \FK(x^*K_\delta x)=\FK(K_\delta)\,\FK(x)^2\ \ge\ \FK(K_\delta)\,\FK(d)^2 .
\]
Take $d=\FK(R)^{1/2}R^{-1/2}$, which lies in $R(H)$ because $R\ge\delta I$, and has $\FK(d)=1$ and
$\tr(d^*Rd)=\FK(R)$. The two displays give $\FK(R)\ge\FK(K_\delta)$. Only this inequality is used; equality, which is
Szeg\H{o}'s formula in the setting of \cite{Lab05}, is not needed.

\emph{Step 4: finite windows.} Let $W\subseteq\Ga$ be finite, and let $g_1H>g_2H>\dots>g_mH$ be the cosets that meet
$W$, with blocks $W_j=W\cap g_jH$. Reveal $\omega$ under $\mu_{K_\delta}$ block by block, $W_1$ first. Fix $j$ and any
values on $P=W_1\cup\dots\cup W_{j-1}$. The restriction of $\mu_{K_\delta}$ to $P\cup W_j$ is determinantal with kernel
the compression of $K_\delta$, which lies between $\delta I$ and $(1-\delta)I$. Lemma~\ref{lem:cond}, with
$\varepsilon=\delta$ and $A\cup B=P$, shows that the conditional law on $W_j$ is determinantal with a kernel $K'$ satisfying
\[
K'\ \ge\ \bigl((K_\delta)_{P\cup W_j}\bigr)^{\langle P\rangle}\ \ge\ \bigl(K_\delta^{\langle U(g_j)\rangle}\bigr)_{W_j},
\qquad U(g_j)=\{x:\ xH>g_jH\},
\]
where the middle term is the Schur complement relative to $P$ of the compression of $K_\delta$ to $P\cup W_j$, which is
what Lemma~\ref{lem:cond} gives, and the right side is defined because the compression of $K_\delta$ to
$\ltwo(U(g_j))$ is at least $\delta$. The second inequality holds because $P\subseteq U(g_j)$, so that more vectors $c$ are admitted in the infimum. Left
multiplication by $g_j^{-1}$ maps $g_jH$ onto $H$ and $U(g_j)$ onto $U_+$, by left invariance, and $\lambda(g_j)$
commutes with $K_\delta$; so the right side is unitarily equivalent, through $\lambda(g_j)$, to the compression of $R$
to $g_j^{-1}W_j\subseteq H$. By monotonicity of determinantal measures in the kernel, which needs no commutation
\cite[Thm.~2.9]{LyICM}, the conditional law on $W_j$ dominates the translate of $\mu_R$ restricted to $g_j^{-1}W_j$. By
the hypothesis on $H$ and Step~3, $\pi_{\FK(R)}\dom\mu_R$ and $\FK(R)\ge\FK(K_\delta)$. So, whatever was revealed
before, the conditional law on each block dominates $\pi_{\FK(K_\delta)}$ on that block. Coupling block by block, with
Strassen's theorem \cite{Str65} on each block and the blocks composed sequentially as in \cite[Lemma~5.2]{LS03}, gives $\pi_{\FK(K_\delta)}\dom\mu_{K_\delta}$ on $W$, and letting $W$ increase gives it on $\Ga$.

\emph{Step 5: limits.} Since $K_\delta\ge(1-2\delta)K$ and $\FK$ is monotone, $\FK(K_\delta)\ge(1-2\delta)\FK(K)$, so
$\pi_{(1-2\delta)\FK(K)}\dom\mu_{K_\delta}$. As $\delta\to0$, $K_\delta\to K$ in norm, so every cylinder probability
$\det((K_\delta)_F)$ converges to $\det(K_F)$; domination passes to the weak limit and gives $\pi_{\FK(K)}\dom\mu_K$.
Applied to $I-K$, which is also equivariant, this gives $\pi_{\FK(I-K)}\dom\mu_{I-K}$, which is
$\mu_K\dom\pi_{1-\FK(I-K)}$ as in the proof that (2) and (3) of Theorem~\ref{thm:A} are equivalent.
\end{proof}

\begin{proof}[Proof of Corollaries~\ref{cor:lo} and~\ref{cor:finite}]
For Corollary~\ref{cor:lo} take $H=\{e\}$: an equivariant $R$ on $\ltwo(\{e\})$ is a number $r\in[0,1]$, and
$\mu_R=\pi_r=\pi_{\FK(R)}$. For Corollary~\ref{cor:finite} it suffices to check the hypothesis for $H$ of order $2$,
since a left order on the group $\Ga/H$ is an order on the coset space invariant under left multiplication by $\Ga$.
An $H$-equivariant $0\le R\le I$ on $\ltwo(H)$ has the form $R=\begin{pmatrix}a&b\\b&a\end{pmatrix}$ with $a,b$ real and
$0\le a-|b|\le a+|b|\le1$. With the normalised trace, $p:=\FK(R)=\sqrt{(a+b)(a-b)}=\sqrt{a^2-b^2}$. On two points
stochastic domination amounts to the inequalities for the four nontrivial increasing events (a given point is present,
both are, at least one is) \cite{Str65}. Each single point lies
in $\omega$ with probability $a\ge p$; both lie in $\omega$ with probability $\det R=a^2-b^2=p^2$; and at least one lies
in $\omega$ with probability $1-\det(I-R)=1-(1-a)^2+b^2\ge1-(1-p)^2$, because $0\le1-a\le1-p$. So
$\pi_{\FK(R)}\dom\mu_R$.
\end{proof}

\begin{remark}\label{rem:block}
The subgroup $H$ has to be treated as a block. Let $H$ be finite of order $n$, let $R$ be $H$-equivariant, and reveal
the points of $H$ in any order $x_1,\dots,x_n$. The conditional probabilities that $x_j$ is present given that
$x_1,\dots,x_{j-1}$ are present multiply to $\det R=\FK(R)^n$, so either one of them is smaller than $\FK(R)$ or all of
them equal $\FK(R)$. For $n=2$ and $R$ as in the proof above with $0<|b|<a$ they are $a$ and $p^2/a<p$. So the criterion that
compares each conditional probability with $\FK(R)$ fails along every revelation order of $H$, although
$\pi_{\FK(R)}\dom\mu_R$ holds. For the same reason the invariant random
orders that Austin uses in place of an invariant order \cite[Thm.~B]{Aus} do not by themselves give the lower
inequality on every group: his determinant of the absolutely continuous part is a Fuglede--Kadison determinant
\cite[Prop.~2.8, Def.~2.11]{Aus}, equal to $\FK(K)$ for the functional of $K\in R(\Ga)$, and its logarithm is the mean, over the random order, of the logarithm
of the Schur complement at the identity relative to the past, and a mean gives no bound at each site. Since a finite set
with a total order invariant under a transitive action is a single point, Theorem~\ref{thm:B} gives nothing for a finite
group beyond its hypothesis.
\end{remark}

\section*{Acknowledgments and funding}

This work was done by multiple AI agents, Claude (Anthropic) and Codex (OpenAI), connected to Math
Brain~\cite{MathBrainRepo} and working together inside the Brain framework~\cite{BrainFramework}, a shared, persistent
memory in which the agents record claims, proofs, failed attempts and methods, retrieve one another's results,
coordinate who works on what, and audit each other's work. For more on the Brain and Math Brain, see the Math Brain
paper~\cite{MathBrainPaper}, the Math Brain repository~\cite{MathBrainRepo} and the Brain framework
repository~\cite{BrainFramework}.

The research presented in this paper was supported by the European Research Council (ERC) under the European
Union's Horizon Europe research and innovation programme (grant agreement No.\ 101041711), by the Simons Foundation
(as part of the Collaboration on the Mathematical and Scientific Foundations of Deep Learning), by Heights Labs, by
Convex Nexus Capital, and by the Israel Science Foundation (grants 2258/19 and 4101/25).

\section*{Code and data availability}

The paper uses no data. The two Python scripts in the directory \texttt{checks} of the paper's repository check
Lemma~\ref{lem:cond} numerically, by enumerating all configurations on a set of six points, for every pair of disjoint sets $A$ and $B$ that
leaves a point free and for three random kernels.

\section*{Declaration of competing interest}

The author holds equity in Heights Labs and Convex Nexus Capital and has a management role at Convex Nexus Capital.

\section*{Declaration of generative AI and AI-assisted technologies}

This paper is part of a broader research project on proving open problems with AI systems. The arguments and the
computer checks were developed with AI systems (OpenAI Codex and Anthropic Claude agents), which also drafted the
text. The conjecture is due to Lyons \cite{LyICM} and the case of $\mathbb Z^d$ to Lyons and Steif \cite{LS03}. The
proof of Theorem~\ref{thm:A} rests on the approximation theorem for the Fuglede--Kadison determinant of Li and Thom
\cite{LT14}, and the proof of Theorem~\ref{thm:B} on the Jensen inequality for subdiagonal algebras, due to Labuschagne
\cite{Lab05} and Blecher and Labuschagne \cite{BL}. The author directed the work and is responsible for the content of the paper and
for decisions concerning its dissemination.

\appendix
\section{Notation}\label{app:notation}

Symbols used in more than one place, in the order in which they first appear. Letters used only inside one proof are
defined where they are used.

\begin{center}\begin{tabular}{@{}lp{0.56\textwidth}l@{}}\toprule
Symbol & Meaning & First appears \\ \midrule
$\Ga$ & a countable group & Section~\ref{sec:intro} \\
$K$ & a $\Ga$-equivariant operator on $\ltwo(\Ga)$ with $0\le K\le I$ & Section~\ref{sec:intro} \\
$I$ & the identity operator & Section~\ref{sec:intro} \\
$\mu_K$ & the determinantal measure on $\{0,1\}^{\Ga}$ given by $K$ & Section~\ref{sec:intro} \\
$\pi_p$ & the product measure on $\{0,1\}^{\Ga}$ of density $p$ & Section~\ref{sec:intro} \\
$\dom$ & stochastic domination & Section~\ref{sec:intro} \\
$\FK(T)$ & the Fuglede--Kadison determinant $\exp\tr(\log T)$ & Section~\ref{sec:intro} \\
$H$ & a subgroup of $\Ga$ with $\Ga/H$ totally ordered & Section~\ref{sec:intro} \\
$R$ & an $H$-equivariant operator on $\ltwo(H)$ with $0\le R\le I$ & Theorem~\ref{thm:B} \\
$C_2$, $F_2$ & the group of order two; the free group on two generators & Section~\ref{sec:intro} \\
$e_x$ & the standard basis vector of $\ltwo(\Ga)$ at $x\in\Ga$ & Section~\ref{sec:notation} \\
$\lambda$, $\rho$ & the left and right regular representations of $\Ga$ & Section~\ref{sec:notation} \\
$R(\Ga)$ & the von Neumann algebra generated by $\rho(\Ga)$ & Section~\ref{sec:notation} \\
$\tr$ & the trace $\tr(T)=\langle Te_e,e_e\rangle$ on $R(\Ga)$ & Section~\ref{sec:notation} \\
$T^*$ & the adjoint of an operator $T$ & Section~\ref{sec:notation} \\
$\omega$ & a configuration in $\{0,1\}^{\Ga}$, identified with a subset of $\Ga$ & Section~\ref{sec:notation} \\
$U_+$, $S$ & $\{x:xH>H\}$ and $U_+^{-1}$ & Section~\ref{sec:B} \\
$\mathcal A$, $\mathcal A_0$ & the subdiagonal algebra and its ideal & Section~\ref{sec:B} \\
$\Phi$ & the conditional expectation of $R(\Ga)$ onto $R(H)$ & Section~\ref{sec:B} \\
$K_\delta$ & $(1-2\delta)K+\delta I$ & Section~\ref{sec:B} \\
\bottomrule\end{tabular}\end{center}

\begin{thebibliography}{99}\raggedright
\bibitem{Arv67} W. Arveson, Analyticity in operator algebras, \emph{Amer. J. Math.} 89 (1967), 578--642.
\bibitem{Aus} T. Austin, Entropy and determinants for unitary representations, arXiv:2412.13751v4, 2026.
\bibitem{BL} D. P. Blecher and L. E. Labuschagne, Applications of the Fuglede--Kadison determinant: Szeg\H{o}'s
theorem and outers for noncommutative $H^p$, \emph{Trans. Amer. Math. Soc.} 360 (2008), no.~11, 6131--6147; arXiv:math/0609662 (section and corollary numbers
cited here are those of the arXiv version).
\bibitem{Cyr21} J. Cyr, Stationary determinantal processes on $\mathbb Z^d$ with $N$ labeled objects per site, Part I:
basic properties and full domination, \emph{J. Theoret. Probab.} 34 (2021), no.~3, 1321--1365,
doi:10.1007/s10959-020-01062-5.
\bibitem{Kes59} H. Kesten, Symmetric random walks on groups, \emph{Trans. Amer. Math. Soc.} 92 (1959), 336--354.
\bibitem{Lab05} L. E. Labuschagne, A noncommutative Szeg\H{o} theorem for subdiagonal subalgebras of von Neumann
algebras, \emph{Proc. Amer. Math. Soc.} 133 (2005), 3643--3646.
\bibitem{LT14} H. Li and A. Thom, Entropy, determinants, and $L^2$-torsion, \emph{J. Amer. Math. Soc.} 27 (2014),
239--292.
\bibitem{Ly03} R. Lyons, Determinantal probability measures, \emph{Publ. Math. Inst. Hautes \'Etudes Sci.} 98 (2003),
167--212.
\bibitem{LyT16} R. Lyons and A. Thom, Invariant coupling of determinantal measures on sofic groups, \emph{Ergodic
Theory Dynam. Systems} 36 (2016), 574--607; arXiv:1402.0969.
\bibitem{LyICM} R. Lyons, Determinantal probability: basic properties and conjectures, \emph{Proc. Int. Congress
Math.} (Seoul, 2014), vol. IV, 137--161; arXiv:1406.2707.
\bibitem{LS03} R. Lyons and J. E. Steif, Stationary determinantal processes: phase multiplicity, Bernoullicity,
entropy, and domination, \emph{Duke Math. J.} 120 (2003), 515--575.
\bibitem{RS} G. Ray and Y. Spinka, Characterizations of amenability through stochastic domination and finitary codings,
\emph{Trans. Amer. Math. Soc.} (2026), doi:10.1090/tran/9703; arXiv:2304.13784.
\bibitem{Str65} V. Strassen, The existence of probability measures with given marginals, \emph{Ann. Math. Statist.}
36 (1965), 423--439.
\bibitem{MathBrainPaper}
Y.~Shmalo, \emph{Math Brain: shared, auditable memory for multi-agent mathematical research}, October 2026.
\url{https://github.com/yspennstate/math-brain-public/tree/5427734deee9be643d6dfbdd5e7d4baeeefe4a4a/papers/math-brain-2026-10-24}
\bibitem{MathBrainRepo}
Y.~Shmalo, \emph{Math Brain: public snapshot, code and tests}, 2026.
\url{https://github.com/yspennstate/math-brain-public/tree/5427734deee9be643d6dfbdd5e7d4baeeefe4a4a}
\bibitem{BrainFramework}
Y.~Shmalo, \emph{Brain Framework}, public software repository, 2026.
\url{https://github.com/yspennstate/brain-framework-public/tree/46b935d26d5a2893f7790b60466387dd8da47ef5}
\end{thebibliography}

\end{document}
